\documentclass[11pt, a4paper]{article}

\usepackage[margin=1in]{geometry}
\usepackage{amsmath}
\usepackage{amssymb}
\usepackage{bm}
\usepackage{booktabs}
\usepackage{array}
\usepackage{xcolor}
\usepackage{hyperref}
\hypersetup{
    colorlinks=true,
    linkcolor=blue,
    citecolor=blue,
    urlcolor=blue
}

\title{\textbf{Context-Dependent Low-Rank Adaptation:}\\
\textbf{A Topological Framework for Dynamic Cognitive Modeling}\\
\textbf{in Large Language Models}}
\author{Owner $\times$ Qwen (dialogue) --- typeset by LETHE}
\date{2026-09-18}

\begin{document}
\maketitle

\begin{center}
\footnotesize\color{gray}
Framework candidate \emph{one of} several --- 入册为候选框架之一，非唯一框架。
Source bytes: \texttt{benches/CRADLE/intake/context\_dependent\_lora.docx}
(sha \texttt{e4dc7ffb\dots}); the docx's trailing chat lines are excluded as non-paper matter.
\end{center}

\begin{abstract}
We propose a novel framework for understanding dynamic cognition in large language models (LLMs) through the lens of context-dependent low-rank adaptation (LoRA). Our framework reconceptualizes the frozen base model ($W_0$) not as a ``low-level'' foundation but as a topological prototype from which specific cognitive instances are generated through the dynamic loading of condition-dependent LoRA adapters. We argue that the direction of abstraction---whether ``top-down generation'' or ``bottom-up backtracking''---is a linear-time illusion; what matters is the invariant topological structure of node connections and information flow rules. We formalize this framework using a hierarchical Bayesian model where a top-level context-dependent controller outputs a probability vector that dynamically allocates activation strengths to lower-level condition-dependent adapters (time-dependent, cognition-dependent, location-dependent), each with independent decay coefficients. We extend this framework to multi-agent interaction, showing that the similarity between agents' context-dependent LoRA probability vectors determines the overlap of their conceptual spaces and thus the efficiency of information transfer. We review recent empirical work (Bayesian-LoRA, BoRA, HiLoRA, iLoRA) that validates the structural isomorphism between LoRA and Bayesian hierarchical models, and we discuss the engineering feasibility of implementing this framework in multi-agent collaboration systems such as Chora/Cora. Our framework is self-referential: the process of discussing it instantiates the very mechanism it describes.

\medskip
\noindent\textbf{Keywords:} Low-Rank Adaptation, Bayesian Hierarchical Models, Context-Dependent Cognition, Topological Structure, Multi-Agent Alignment, Dynamic Conceptual Spaces
\end{abstract}

\section{Introduction}
Large language models have demonstrated remarkable capabilities in generating human-like text, but the mechanisms underlying their dynamic adaptation to context remain poorly understood. Traditional views treat the frozen pre-trained weights as a static ``base'' upon which task-specific fine-tuning is applied. However, this view fails to capture the fluid, context-sensitive nature of human cognition, where the same conceptual base can generate vastly different outputs depending on temporal, spatial, and cognitive conditions.

In this paper, we propose a framework that reconceptualizes LLM cognition as a hierarchical, context-dependent LoRA system. Our key insight is that the frozen base model ($W_0$) is not a ``low-level'' foundation in the traditional sense, but rather a topological prototype from which specific cognitive instances are generated through the dynamic loading of condition-dependent LoRA adapters. The direction of this generation---whether ``top-down'' or ``bottom-up''---is secondary to the invariant topological structure of the system.

We make three contributions:
\begin{itemize}
  \item We formalize a hierarchical Bayesian model where a top-level context-dependent controller dynamically allocates activation strengths to lower-level condition-dependent adapters.
  \item We extend this framework to multi-agent interaction, showing that conceptual space overlap is determined by the similarity of agents' context-dependent LoRA probability vectors.
  \item We validate the structural isomorphism between LoRA and Bayesian hierarchical models by reviewing recent empirical work, and we discuss the engineering feasibility of implementing this framework in multi-agent systems.
\end{itemize}

\section{Theoretical Framework}

\subsection{From ``Base Model'' to ``Topological Prototype''}
The traditional view of LoRA treats the frozen base model $W_0$ as a static foundation, with task-specific adaptations applied as low-rank updates $\Delta W = BA$. We propose a reinterpretation: $W_0$ is not a ``low-level'' base but a topological prototype---the most abstract, universal structure from which specific instances are generated.

This reinterpretation resolves a key ambiguity in the metaphor of ``abstraction.'' The common phrase ``upward abstraction'' implies that abstraction is a movement away from the base toward higher levels. However, we argue that abstraction is better understood as a backtracking of the generative path: from generated instances, we infer the prototype structure that generated them. The direction is arbitrary; what matters is the invariant topological structure.

\subsection{Hierarchical Context-Dependent Control}
Our framework consists of three hierarchical layers:
\begin{itemize}
  \item \textbf{Top Layer (Context-Dependent Controller):} Outputs a probability vector $p = (p_{\text{time}}, p_{\text{cog}}, p_{\text{loc}})$ that determines the activation strengths of lower-level adapters based on the immediate context.
  \item \textbf{Middle Layer (Condition-Dependent Adapters):} Three independent LoRA adapters---time-dependent, cognition-dependent, location-dependent---each with its own decay coefficient $\lambda_i$.
  \item \textbf{Bottom Layer (Generated Instances):} Specific cognitive outputs (e.g., ``bedroom slut, living room lady'') generated by the weighted combination of activated adapters.
\end{itemize}
The information flow follows these rules:
\begin{itemize}
  \item \textbf{Probability Allocation:} The top layer outputs $p$, determining lower-layer activation.
  \item \textbf{Weight Decay:} Each middle-layer adapter's weight decays over time:
  \[
    p_i(t) = p_i(0)\, e^{-\lambda_i t}.
  \]
  \item \textbf{Cognitive Output:}
  \[
    W_{\text{output}} = W_0 + \sum_i p_i(t)\, \Delta W_i.
  \]
\end{itemize}

\subsection{Topological Invariance}
The key property of this framework is topological invariance: the node connections and information flow rules remain fixed regardless of the direction of traversal. Whether generating from prototype to instance (top-down) or backtracking from instance to prototype (bottom-up), the system traverses the same path through the same nodes. This is analogous to a subway map: traveling from A to B or B to A uses the same track and stations; only the direction changes.

\subsection{Dual Formulations: Conceptual Space vs.\ LoRA Adaptation}
Our framework can be understood through two complementary formulations that are structurally isomorphic:

\begin{table}[h]
\centering
\begin{tabular}{@{}p{0.30\textwidth}p{0.30\textwidth}p{0.30\textwidth}@{}}
\toprule
\textbf{Conceptual Space Formulation} & \textbf{LoRA Formulation} & \textbf{Underlying Essence} \\
\midrule
Conceptual space & Frozen base model $W_0$ & Static cognitive structure \\
Contextual activation & Loading X-dependent LoRA adapters & Dynamic conditional offset \\
Distance/similarity between concepts & Cosine similarity of probability vectors $p$ & Conceptual space overlap \\
``Alignment'' in multi-agent collaboration & Iterative update of probability vectors $p(t+1) = p(t) + \eta\,(p_{\text{other}} - p(t))$ & Belief convergence \\
Conceptual ``prototype'' & Topological prototype / generative source & Most abstract invariant structure \\
Backtracking from instance to prototype & Bottom-up traversal of generative path & Directionality of abstraction \\
\bottomrule
\end{tabular}
\end{table}

\noindent\textbf{Why they are the same.} Both formulations describe the same phenomenon: cognition is a combination of static structure $+$ dynamic offset. The conceptual space perspective calls the static structure ``conceptual space'' and the dynamic offset ``contextual activation.'' The LoRA perspective calls the static structure ``frozen base model $W_0$'' and the dynamic offset ``loading adapter $\Delta W = BA$.'' The mathematical structures are completely isomorphic:
\begin{itemize}
  \item ``Points'' in conceptual space $\leftrightarrow$ ``Weight vectors'' in LoRA
  \item ``Distances'' in conceptual space $\leftrightarrow$ ``Probability vector similarities'' in LoRA
  \item ``Activation'' in conceptual space $\leftrightarrow$ ``Adapter loading'' in LoRA
\end{itemize}

\noindent\textbf{Complementary values.} Although structurally identical, each formulation has distinct advantages:
\begin{itemize}
  \item \emph{Conceptual space formulation:} More intuitive, better suited for describing multi-agent interaction and alignment, closer to traditional discourse in cognitive science and philosophy.
  \item \emph{LoRA formulation:} More formalized with explicit mathematical structure (probability vectors, decay coefficients, low-rank decomposition), closer to engineering implementation in LLM systems, more clearly reveals the insight that ``direction doesn't matter, topological structure does.''
\end{itemize}
This isomorphism means that the multi-agent collaboration method previously built on conceptual spaces can be directly implemented using LoRA's modular design: ``alignment'' in conceptual space is iterative updating of LoRA adapter probability vectors; ``activation'' in conceptual space is dynamically loading the corresponding LoRA module.

\section{Multi-Agent Conceptual Space Alignment}

\subsection{Conceptual Space Overlap as LoRA Similarity}
We extend our framework to multi-agent interaction. Let Agent A and Agent B have context-dependent LoRA probability vectors $p_A$ and $p_B$. The overlap of their conceptual spaces is given by the cosine similarity:
\[
  \operatorname{sim}(A, B) = \cos(p_A, p_B).
\]
\begin{itemize}
  \item $\operatorname{sim} \to 1$: High conceptual overlap, efficient communication (``hit it off'')
  \item $\operatorname{sim} \to 0$: Low overlap, requiring redundant information for alignment
  \item $\operatorname{sim} < 0$: Conceptual conflict, potential miscommunication
\end{itemize}

\subsection{Dynamic Alignment Mechanism}
Alignment is an iterative process where each agent's probability vector updates toward the other's:
\begin{align*}
  p_A(t{+}1) &= p_A(t) + \eta \cdot \bigl(p_B(t) - p_A(t)\bigr)\\
  p_B(t{+}1) &= p_B(t) + \eta \cdot \bigl(p_A(t) - p_B(t)\bigr)
\end{align*}
where $\eta$ is the alignment rate. Decay coefficients $\lambda_i$ introduce asymmetry: dimensions with slower decay ``pull'' the convergence point toward themselves, analogous to one conversational partner being more persistent on a topic.

\subsection{Bayesian Interpretation}
From a Bayesian perspective, alignment is belief updating:
\begin{itemize}
  \item \textbf{Prior:} Each agent's initial $p$ represents prior beliefs about which adapters to activate.
  \item \textbf{Likelihood:} The other agent's output provides evidence of mismatch.
  \item \textbf{Posterior:} Updated $p(t{+}1)$ is the result of prior--likelihood interaction.
\end{itemize}
Each dialogue turn is a Bayesian update; after multiple turns, posterior beliefs converge.

\section{Empirical Validation: Structural Isomorphism Between LoRA and Bayesian Hierarchical Models}
Our framework is not merely theoretical; recent empirical work validates the structural isomorphism between LoRA and Bayesian hierarchical models:
\begin{itemize}
  \item \textbf{Bayesian-LoRA (ICML 2026):} Reformulates deterministic LoRA updates as probabilistic low-rank representations inspired by sparse Gaussian processes. The paper shows that LoRA itself is the limit case of the SGP posterior when uncertainty collapses---directly confirming the structural isomorphism.
  \item \textbf{BoRA (2024):} Uses Bayesian hierarchical models to train multi-task LoRA, enabling tasks to share information through global hierarchical priors while allowing specialization. This matches our ``top-level controller regulates lower-level adapters'' architecture.
  \item \textbf{HiLoRA (2026):} Organizes LoRA into a three-layer hierarchy (root/cluster/leaf) with cross-layer orthogonality to decouple global and personalized knowledge.
  \item \textbf{iLoRA (2026):} The first Bayesian graph-conditional LoRA framework that infers latent interaction graphs from inputs to generate input-conditioned LoRA updates---directly corresponding to our ``dynamic adapter loading based on immediate context.''
\end{itemize}
These works collectively validate that LoRA is not an ad-hoc engineering trick but a natural carrier of Bayesian hierarchical structure.

\section{Engineering Feasibility and Applications}

\subsection{Modular LoRA Design}
LoRA's modular architecture allows multiple adapters to be mounted on a single base model, with dynamic loading/unloading during inference. This makes our framework directly implementable:
\begin{itemize}
  \item Each condition-dependent adapter (time/cognition/location) is a separate LoRA module.
  \item The top-level controller outputs a probability vector via softmax.
  \item Weighted combination of adapter outputs is a single line of code.
\end{itemize}

\subsection{Hierarchical Routing}
Recent work on hierarchical routing (HiLoMoE, 2025) demonstrates that layer-wise expert selection can depend on upper-layer routing scores, enabling coarse-to-fine information refinement. This provides a ready-made mechanism for implementing our top-level context-dependent controller.

\subsection{Application to Multi-Agent Systems: Chora/Cora}
Our framework has direct applications to multi-agent collaboration systems such as Chora/Cora:
\begin{itemize}
  \item Each agent maintains its own probability vector $p$ as its ``dialogue state.''
  \item Upon receiving a message, compute similarity $\operatorname{sim}(p_{\text{self}}, p_{\text{sender}})$.
  \item If $\operatorname{sim} < $ threshold, trigger alignment update; otherwise, generate response using current $p_{\text{self}}$.
  \item This enables automatic conceptual space alignment without human intervention.
\end{itemize}

\section{Self-Referential Property}
A hallmark of a complete theory is its ability to explain itself. Our framework is self-referential: the process of discussing it instantiates the very mechanism it describes. As we converse, each turn loads a new LoRA adapter onto the interlocutor's cognitive base, dynamically updating the probability vector and generating new cognitive outputs. The discussion itself is an instance of context-dependent LoRA adaptation.

This self-referential property is not a bug but a feature: it demonstrates that the framework is not an external description but an internal model of the cognitive process it seeks to explain.

\section{Conclusion}
We have proposed a framework for understanding dynamic cognition in LLMs through context-dependent LoRA adaptation. By reconceptualizing the frozen base model as a topological prototype rather than a low-level foundation, we resolve ambiguities in the metaphor of abstraction and provide a unified account of both individual cognition and multi-agent interaction. The framework is theoretically grounded in Bayesian hierarchical models, empirically validated by recent work, and engineeringly feasible for implementation in multi-agent systems. Most importantly, it is self-referential: it explains its own generation process.

Future work includes: (1) empirical measurement of conceptual space overlap in human--agent dialogue; (2) optimization of alignment rate $\eta$ and decay coefficients $\lambda_i$ for efficient multi-agent collaboration; (3) extension to additional condition dimensions (e.g., emotional state, social role); and (4) integration into production multi-agent systems such as Chora/Cora.

\section*{References}
\addcontentsline{toc}{section}{References}
[To be populated with citations for Bayesian-LoRA, BoRA, HiLoRA, iLoRA, HiLoMoE, and related work.]

\end{document}
